\documentclass[a4paper, noanswers, addpoints, 12pt]{exam}
\usepackage[T1]{fontenc}
\usepackage[utf8]{inputenc}
\usepackage{lmodern}
\usepackage{tikz}
\usetikzlibrary{positioning,shapes.multipart,arrows.meta,shadows}
\usepackage[french]{babel}
\usepackage{xcolor}
\usepackage{listings}
\usepackage{graphicx}
\usepackage[most]{tcolorbox}
\usepackage{minted}

\definecolor{codegreen}{rgb}{0,0.6,0}
\definecolor{codegray}{rgb}{0.5,0.5,0.5}
\definecolor{codepurple}{rgb}{0.58,0,0.82}
\definecolor{backcolour}{rgb}{0.95,0.95,0.92}
 
\lstdefinestyle{mystyle}{
    backgroundcolor=\color{backcolour},   
    commentstyle=\color{codegreen},
    keywordstyle=\color{magenta},
    numberstyle=\tiny\color{codegray},
    stringstyle=\color{codepurple},
    basicstyle=\footnotesize,
    breakatwhitespace=false,         
    breaklines=true,                 
    captionpos=b,                    
    keepspaces=true,                 
    numbers=left,                    
    numbersep=5pt,                  
    showspaces=false,                
    showstringspaces=false,
    showtabs=false,                  
    tabsize=2
}
\lstset{style=mystyle, language=Python}

\title{CC Term NSI - Arbre et graphes}
\author{Babash - NSI Tani Malandi et Johanne Bakalara}
\date{Décembre 2025, Durée: 2h}


\usepackage[a4paper,top=2.3cm,bottom=2.8cm,left=2cm,right=2cm,headsep=0.6cm,footskip=1.2cm]{geometry}
\begin{document}
\runningfooter{}{\footnotesize Babash - NSI Tani Malandi et Johanne Bakalara \textperiodcentered{} CC BY-SA 4.0 \textperiodcentered{} page \thepage}{}
\firstpagefooter{}{\footnotesize Babash - NSI Tani Malandi et Johanne Bakalara \textperiodcentered{} CC BY-SA 4.0 \textperiodcentered{} page \thepage}{}

\maketitle
\thispagestyle{headandfoot}
\begin{center}\footnotesize\itshape Exercice sur les arbres d'après l'exercice 1 du sujet de NSI du baccalauréat 2023 (jour 2, Polynésie, code 23-NSIJ2PO1).\end{center}

\pagestyle{headandfoot}
\thispagestyle{headandfoot}
\firstpageheader{\makebox[\textwidth]{Nom : \dotfill\ Pr\'enom : \dotfill\ Classe : \makebox[3cm]{\dotfill}}}{}{}


\section*{Consignes:}
\begin{center} 
\fbox{\parbox{\linewidth}{
\begin{center}
Rédigez vos réponses sur une feuille libre.\\
La clarté de la rédaction et les explications font partie intégrante de la notation. Une réponse non rédigée ou sans explication sera considérée fausse.\\
Encadrez clairement vos réponses pour une meilleure lisibilité.
\end{center}


}}
\end{center}

\vsword{Note}
\hsword{Note}

\section*{Barème :}
\begin{center}
  \par\bigskip
  \gradetable[h]
\end{center}
\begin{questions}

\section{Notions fondamentales sur les arbres : }

\question Question de cours :
\begin{parts}
\part[1] Recopier et compléter le texte ci-dessous en choisissant des expressions parmi
au maximum, au minimum, exactement, feuille, racine, sous-arbre
gauche et sous-arbre droit :
\begin{parts}
    \item Le noeud initial est appelé ... . 
    \item Un noeud qui n'a pas d'enfants est appelé ....
    \item Un arbre binaire est un arbre dans lequel chaque noeud a ... deux enfants.
    \item Un arbre binaire de recherche est un arbre binaire dans lequel tout
    noeud est associé à une clé qui est :
    \begin{itemize}
        \item supérieure à chaque clé de tous les noeuds de son ...
        \item inférieure à chaque clé de tous les noeuds de son ....
    \end{itemize}
\end{parts}

\newpage

\begin{figure}
    \centering
    {\centering
    \begin{tikzpicture}[level distance=1.4cm,every node/.style={draw,circle,minimum size=0.8cm,inner sep=1pt},
 level 1/.style={sibling distance=3.2cm},level 2/.style={sibling distance=2cm}]
\node {13} child {node {6} child {node {7}} child {node {5} child[missing] child {node {1}}}} child {node {4}};
\end{tikzpicture}

    \par}
    \caption{Arbre 1}
    \label{fig:arbre1}
\end{figure}

\begin{figure}
    \centering
    {\centering
    \begin{tikzpicture}[level distance=1.4cm,every node/.style={draw,circle,minimum size=0.8cm,inner sep=1pt},
 level 1/.style={sibling distance=3.4cm},level 2/.style={sibling distance=1.6cm}]
\node {6} child {node {4} child {node {3}} child {node {5}}} child {node {8} child {node {7}} child {node {9}}};
\end{tikzpicture}

    \par}
    \caption{Arbre 2}
    \label{fig:arbre2}
\end{figure}

\setcounter{partno}{1}
\part[3] Pour chacun des deux arbres (\ref{fig:arbre1} et \ref{fig:arbre2})
\begin{parts}
    \item Déterminez si chaque arbre représenté est un arbre binaire. Justifiez votre réponse en expliquant pourquoi.
    \item Vérifiez si chaque arbre est un arbre binaire de recherche (ABR) en expliquant pourquoi il satisfait (ou non) les propriétés d'un ABR.
    \item Donner la hauteur de l'arbre
    \item Donner le résultat d'un parcours préfixe sur l'arbre.
    \item Donner le résultat d'un parcours infixe sur l'arbre.
    \item Donner le résultat d'un parcours suffixe sur l'arbre.
\end{parts}

\end{parts}

\vspace{1cm}

\section{Programmation et implémentation des arbres : }

On considère la classe \textit{Noeud}, nous permettant de définir les arbres binaires, définie de la manière suivante en {\tt python} :

\begin{lstlisting}[style=mystyle, language=Python]
class Noeud:
    def __init__(self, v, g=None, d=None):
        """cree un noeud d'un arbre binaire"""
        self.gauche = g
        self.droit = d
        self.valeur = v
\end{lstlisting}

Grâce à cette structure de données, on peut implémenter l'arbre de la figure~\ref{fig:arbre1} de la manière suivante.
\begin{lstlisting}[style=mystyle, language=Python]
n4 = Noeud(4)
n7 = Noeud(7)
n1 = Noeud(1)
n5 = Noeud(5, None, n1)
n6 = Noeud(6, n7, n5)
arbre1 = Noeud(13, n6, n4)
\end{lstlisting}

\question[1] Donner le code permettant d'implémenter l'arbre~\ref{fig:arbre2}.

\question[1] Manipulation des objets :
\begin{parts}
    \item Quel est le résultat de l'expression~\textit{arbre1.gauche.droit.valeur} ?
    \item Donner la commande permettant d'accéder à la valeur du noeud 1 depuis la racine.
\end{parts}

\section{Manipulation récursive des arbres : }

% fonction est feuille à compléter
\question[1] Compléter la fonction~\textit{est\_feuille} ci-dessous permettant de tester si un objet de type  \textit{Noeud} est une feuille :
\begin{lstlisting}[style=mystyle, language=Python]
def est_feuille(abr):
    if ...................................... :
        return ...
    else :
        return False
\end{lstlisting}

\begin{lstlisting}[style=mystyle, language=Python]
def mystere(abr, x):
    if abr is None :
        return 0
    else :
        return 1 + mystere(abr.gauche, x) + mystere(abr.droit, x)
\end{lstlisting}
\newpage
\question[3] Mystère :
\begin{parts}
    \item La fonction \textit{mystere} est-elle récursive ? Justifiez.
    \item Donner la valeur de retour de la fonction \textit{mystere} pour les appels suivants (avec \textit{abr1} et \textit{abr2} sont les arbres des figures~\ref{fig:arbre1} et ~\ref{fig:arbre2}):
    \begin{itemize}
        \item \textit{ mystere(abr1, 9)}
        \item \textit{ mystere(abr1, 5)}
        \item \textit{ mystere(abr2, 6)}
        \item \textit{ mystere(abr2, 7)}
    \end{itemize}
    \item Expliquer ce que fait la fonction \textit{mystere}.
\end{parts}

% pivot
Il existe des algorithmes pour modifier la structure d'un ABR, afin par exemple de diminuer la hauteur d'un ABR, on s'intéresse aux algorithmes appelés rotation, consistant à faire "pivoter" une partie de l'arbre autour d'un de ses noeuds.
L'exemple ci-dessous permet d'expliquer l'algorithme pour réaliser une rotation droite d'un ABR autour de sa racine :

{\centering
    \begin{tikzpicture}[level distance=1.1cm,every node/.style={draw,circle,minimum size=0.7cm,inner sep=1pt,font=\small},
 level 1/.style={sibling distance=2cm},level 2/.style={sibling distance=1.1cm},baseline=(a.base)]
\node (a) at (0,0) {20} child {node {15} child {node {12}} child {node {17}}} child {node {30}};
\draw[->,thick] (2.2,-0.9) -- (3.6,-0.9);
\begin{scope}[xshift=6cm]
\node {15} child {node {12}} child {node {20} child {node {17}} child {node {30}}};
\end{scope}
\end{tikzpicture}

    \par}

\begin{center}
    \textbf{On admet que ces transformations conservent la propriété d'ABR de l'arbre.}
\end{center}

\newpage

La fonction \textit{rotation\_droite} ci-après renvoie une nouvelle instance de type \textit{Noeud}, correspondant à une rotation droite de l'objet de type~\textit{Noeud} à partir duquel elle est appelée :

\begin{lstlisting}[style=mystyle, language=Python]
def rotation_droite(abr):
    pivot = abr.gauche
    abr.gauche = pivot.droit
    pivot.droit = abr
    return Noeud(pivot.valeur, pivot.gauche, pivot.droit)
\end{lstlisting}


Pour réaliser une rotation gauche, on suivra alors l'algorithme suivant :
\begin{itemize}
    \item on appelle pivot le sous-arbre droit de la racine de l'arbre,
    \item le sous-arbre gauche du pivot devient le sous-arbre droit de la racine,
    \item la racine ainsi modifiée devient le sous-arbre gauche du pivot et la racine du
pivot devient la nouvelle racine de l'ABR.
\end{itemize}

\question Compréhension et implémentation de l'équilibrage
\begin{parts}
\part[1] En suivant les différentes étapes de cet algorithme, dessiner l'arbre obtenu
après une rotation droite de l'arbre de la figure~\ref{fig:arbre2}.

\part[1] En suivant les différentes étapes de cet algorithme, dessiner l'arbre obtenu
après une rotation gauche de l'arbre de la figure~\ref{fig:arbre1}.

\part[2] Écrire le code d'une fonction {\tt Python} \textit{rotation\_gauche} qui réalise la
rotation gauche d'un ABR autour de sa racine.
\end{parts}

\section{Graphes : }

Dans cette partie, on considère les villes suivantes : Washington (W) ; Paris (P) ; Berlin
(B) ; Tokyo (T) ; Sydney (S).\\
Les distances des vols entre ces villes, exprimées en milliers de kilomètres, sont les
suivantes :
\begin{itemize}
\item Washington - Paris : 9 ;
\item Washington - Berlin : 4 ;
\item Paris - Berlin : 3 ;
\item Paris - Tokyo : 12 ;
\item Paris - Sydney : 17 ;
\item Berlin - Tokyo : 7 ;
\item Tokyo - Sydney : 7.
\end{itemize}


La compagnie AirInfo propose certains de ces vols :
\begin{itemize}
    \item chaque ville est représentée par un sommet ;
    \item chaque vol direct entre deux villes est représenté par une arête.
\end{itemize}
Le nombre indiqué sur chaque arête est appelé poids : il représente la distance entre
deux villes. \textbf{Cette distance est la même dans un sens ou dans l'autre}.

\question[1] Représenter le graphe des vols de la compagnie AirInfo.\\

On choisit de représenter ce graphe en Python à l'aide d'un dictionnaire dans lequel
chaque clé est un sommet du graphe et la valeur associée à cette clé est elle-même
un dictionnaire, dont les clés sont les villes voisines et les valeurs sont les distances
correspondantes.

\question[1] Compléter le code ci-dessous
\begin{lstlisting}[style=mystyle, language=Python]
 graphe_airinfo = {
    'W': {'P': 9, 'B': 4},
    'P': {...},
    'B': {...},
    ...
    ...: {'P': 17, 'T': 7}
}
\end{lstlisting}
\question[1] Donner la matrice d'adjacence du graphe.
\question[1] Déterminer la valeur de l'expression \textit{graphe\_airinfo['W']['P']}.

\question[1] Écrire le code d'une fonction Python \textit{vol\_direct} permettant de déterminer s'il
existe une liaison directe entre deux villes. Cette fonction prend en paramètres
un dictionnaire graphe représentant le graphe, et deux clés du dictionnaire
\textit{ville1} et \textit{ville2} représentant deux villes, et renvoie \textit{True} si une telle liaison
entre les villes \textit{ville1} et \textit{ville2} existe, c'est-à-dire si les deux sommets
\textit{ville1} et \textit{ville2} sont reliés dans le graphe \textit{graphe} par une arête, \textit{False} sinon.\\

Exemples :\\
$>>>$ \textit{vol\_direct(graphe\_airinfo, 'T', 'B')}\\
\textit{True}\\
$>>>$ \textit{vol\_direct(graphe\_airinfo, 'W', 'B')}\\
\textit{True}
\end{questions}
\end{document}
